\section{Related work}
\label{sec:related}

\paragraph{General-purpose ABM platforms.} NetLogo \citep{wilensky1999netlogo} is widely used in
teaching, with its own language and an integrated interface; it excels at exploratory models,
and a Python analysis reaches it through a bridge such as PyNetLogo
\citep{jaxarozen2018pynetlogo}; GAMA \citep{taillandier2019building} specialises in spatial
models. Repast Simphony \citep{north2013repast} and MASON
\citep{luke2005mason} are mature Java platforms built for performance and large simulations,
and Repast4Py \citep{collier2020repast4py} distributes a single run over MPI processes from
Python; \modelly{} does so too, under a contract that makes the split run
reproduce the single-process one byte for byte (\cref{sec:split}), and runs column phases on
a GPU through CuPy or JAX rather than through a dedicated language (\cref{sec:gpu}). Agents.jl \citep{datseris2024agents} brings a performant, concise framework to Julia,
and FLAME GPU \citep{richmond2023flamegpu} runs agent models on GPUs. Reviews of the field
\citep{railsback2006agent, abar2017agent} have long noted the trade-off between ease of use and speed that
these tools resolve in different ways.

\paragraph{Python.} Mesa \citep{masad2015mesa, kazil2020mesa, terhoeven2025mesa3} provides the agent, model,
space and data-collection abstractions that most Python ABMs use, and a browser
visualisation. AgentPy \citep{foramitti2021agentpy} combines a similar model with experiment
and Sobol tooling built on SALib \citep{herman2017salib}, though its last release (0.1.5) dates from
December 2021; Melodie \citep{yu2023melodie} adds calibration and scenario management with
Cython acceleration; and mesa-frames \citep{mesaframes} represents agent sets as data frames
to vectorise Mesa models. \modelly{} is written from
scratch and shares no code with Mesa, but its programming model is compatible with Mesa's: its
codemod (\cmd{python -m modelly.migrate})
performs the mechanical rewrites of a Mesa model and reports what needs a decision
(schedulers, legacy grids, batch runs); the two Mesa~2.3 models migrated this way, changing only
the framework calls, give bit-identical results on every seed tried, and the migrated Mesa~3.5
models equivalent distributions (\cref{sec:samecode}). Unlike a data-frame representation, \modelly{}'s columns keep
the per-agent object: each agent reads and writes its row as an attribute, so a model can be
vectorised one phase at a time (\cref{sec:storage}).

\paragraph{Calibration, sensitivity analysis and model exploration.} The methods \modelly{}
implements are standard: approximate Bayesian computation \citep{beaumont2002approximate}, with sequential Monte
Carlo \citep{sisson2007sequential, toni2009approximate, beaumont2009adaptive}, already applied to
agent-based models \citep{vandervaart2015calibration, platt2020comparison}, elementary effects
\citep{morris1991factorial}, variance-based indices \citep{sobol2001global, saltelli2010variance,
jansen1999analysis} and the PAWN density-based index \citep{pianosi2015simple,
pianosi2018distribution}. Libraries implement them for generic simulators
\citep{herman2017salib, iwanaga2022salib, klinger2018pyabc, schalte2022pyabc}, and workflow
tools connect them to ABM platforms: the EMA Workbench \citep{kwakkel2017ema}, OpenMOLE
\citep{reuillon2013openmole} and nlrx for NetLogo \citep{salecker2019nlrx}; reviews and tutorials give
guidance on which to use \citep{thiele2014facilitating, tenbroeke2016which, hartig2011statistical,
lee2015complexities}. The same holds for history matching with emulators
\citep{vernon2010galaxy, andrianakis2015bayesian}, the bootstrap particle filter
\citep{gordon1993novel} and the ensemble Kalman filter \citep{evensen2003ensemble}, and
gradient estimators for discrete choices \citep{bengio2013estimating, jang2017categorical}.
\modelly{}'s contribution here is integration rather than new methods:
the samplers share the simulation's seeding and executors, so a calibration or a Sobol analysis
of a model is reproducible, resumable and parallel without extra code.

\paragraph{Documentation and reproducibility.} The ODD protocol \citep{grimm2006odd,
grimm2010odd, grimm2020odd} and its extension to human decisions, ODD+D
\citep{muller2013describing}, are the standard description of an ABM in a publication, TRACE
\citep{grimm2014towards} documents its testing and analysis, and pattern-oriented modelling
\citep{grimm2005pattern} guides its validation. Archiving model code remains uncommon
\citep{janssen2017practice}, and replicating a published model is a study in itself
\citep{axtell1996aligning, wilensky2007making}. \modelly{} generates an ODD skeleton from the
model's code and records each run in a manifest that a reviewer can rerun and verify.

\Cref{tab:features} compares the built-in capabilities of the Python frameworks and of two
widely used platforms in other languages.

\begin{table}[htbp]
  \centering
  \footnotesize
  \caption{Capabilities as documented by each project in October 2026 (Mesa 3.5.1,
  mesa-frames 0.1.0a0, its only release, an alpha, while its development branch
  \code{ef92326} (0.1.1.dev0) is the one measured in \cref{sec:playbooks}, AgentPy 0.1.5, NetLogo 7, Repast4Py 1.3, Agents.jl 7.1). \checkmark{}: built in;
  ext.: through a companion tool that the project or the tool documents for it (EMA Workbench for
  Mesa, ipysimulate for AgentPy, pyNetLogo with SALib for NetLogo \citep{jaxarozen2018pynetlogo},
  EMEWS \citep{ozik2016desktop} for Repast4Py); --: not documented (nlrx's ABC calibration supports
  NetLogo only up to 6.4, and EMEWS documents no sensitivity analysis or calibration for
  Repast4Py). The judgements are the author's, from each project's documentation and source; the
  evidence for every cell is in \code{paper/data/features\_sources.md}. NetLogo's checkpoints are
  \code{export-world} and \code{import-world}; Agents.jl's do not cover its event-queue models.
  Parallel runs in Mesa, AgentPy, NetLogo and Agents.jl are across runs, not one run. Licences
  from the projects' licence files (mesa-frames' development branch; its 0.1.0a0 release is MIT).}
  \label{tab:features}
  \setlength{\tabcolsep}{3pt}
  \begin{tabular}{lccccccc}
    \toprule
    & \modelly{} & Mesa 3 & mesa-frames & AgentPy & NetLogo & Repast4Py & Agents.jl \\
    \midrule
    Language & Python & Python & Python & Python & NetLogo & Python & Julia \\
    Licence & Apache-2.0 & Apache-2.0 & Apache-2.0 & BSD-3 & GPL-2.0+ & BSD-3 & MIT \\
    Per-agent objects & \checkmark & \checkmark & -- & \checkmark & \checkmark & \checkmark & \checkmark \\
    Array or table agent state & \checkmark & -- & \checkmark & -- & -- & -- & -- \\
    Counted cohorts & \checkmark & -- & -- & -- & -- & -- & -- \\
    One run across processes & \checkmark & -- & -- & -- & -- & \checkmark & -- \\
    Parameter sweeps & \checkmark & \checkmark & -- & \checkmark & \checkmark & ext. & \checkmark \\
    Global sensitivity analysis & \checkmark & ext. & -- & \checkmark & ext. & -- & -- \\
    Bayesian calibration & \checkmark & -- & -- & -- & -- & -- & -- \\
    Checkpoint and resume & \checkmark & -- & -- & -- & \checkmark & \checkmark & \checkmark \\
    Run manifests and bundles & \checkmark & -- & -- & -- & -- & -- & -- \\
    Interactive visualisation & \checkmark & \checkmark & -- & ext. & \checkmark & -- & \checkmark \\
    \bottomrule
  \end{tabular}
\end{table}
